设 $k=3$ 或 $4$，取 $(e_0,e_1,e_2,e_3)=(1,\mathbf i,\mathbf j,\mathbf k)$，并记 $z_a=\sum_{\alpha=0}^{k-1}e_\alpha z_{a\alpha}$，其中 $1\le a\le m$。对实的 $C^2$ 函数 $u$，定义
\[
G_a(u)=\sum_\alpha e_\alpha\partial_{a\alpha}u,\qquad
H_{ab}(u)=\sum_{\alpha,\beta}e_\alpha\bar e_\beta\,
       \partial_{a\alpha}\partial_{b\beta}u,\qquad
\partial_{a\alpha}=\partial/\partial z_{a\alpha}.
\]
则 $H(u)$ 为 Hermitian 矩阵，且 $\operatorname{tr}H(u)=\Delta u$。从实 Hessian 到 $H(u)$ 的映射将实秩一矩阵 $\xi\xi^{\mathsf t}$ 映为 $gg^*$，其中 $g_a=\sum_\alpha e_\alpha\xi_{a\alpha}$，并将非负实矩阵映为非负四元数矩阵。对实的 $C^2$ 函数 $u,v$，
\[
H(uv)=uH(v)+vH(u)+G(u)G(v)^*+G(v)G(u)^* .
\]
在四槽情形下，$H$ 在四元数酉坐标变换下具有协变性；在三槽情形下，对所有槽应用相同的关于指标 $a$ 的实正交变换时也具有协变性。若 $u$ 只依赖于 $P\in\mathbb C_{\mathbf k}^m$，取值于 $z_a=x_a+P_a\mathbf i$（若有第四槽，则与第四槽无关），记 $P_a=p_a+\mathbf k q_a$，并采用
\[
\partial_{P_a}=\tfrac12(\partial_{p_a}-\mathbf k\partial_{q_a}),\qquad
\partial_{\bar P_a}=\tfrac12(\partial_{p_a}+\mathbf k\partial_{q_a}).
\]
在这一约定下，
\[
G_a(u)=2(\partial_{\bar P_a}u)\mathbf i,\qquad
H_{ab}(u)=4\partial_{\bar P_a}\partial_{P_b}u .
\]
对于全纯映射 $F=(F_1,\ldots,F_\ell)$（定义于 $P$），其复 Jacobian 为 $J_{\nu a}=\partial_{P_a}F_\nu$ 时，有 $H(\sum_\nu|F_\nu|^2)=4J^*J$。对实光滑函数 $\rho$，若其自变量是复变量 $F_\nu$，则梯度 $G(\rho\circ F)$ 属于 $J^*$ 的四元数像。
